\section{Las hipótesis, una por una}\label{sec:recorrido}
Esta sección reúne en un solo hilo lo que las secciones anteriores presentan por separado. Para cada hipótesis indica qué hecho la motivó, cómo se representó en el modelo, qué resultó y en qué estado queda. La tabla~\ref{tab:recorrido} la resume. Un resultado «compatible» significa que el modelo reproduce lo observado con supuestos razonables; no prueba que el evento haya ocurrido así. El índice exploratorio de compatibilidad (IEC) compara hipótesis entre sí con los mismos criterios; no es una validación.

\subsection{Puntos de entrada del flujo ensayados en el modelo}
\paragraph{H1: el cauce principal sale en el abanico.} \emph{Por qué se planteó.} Las primeras informaciones recogidas por Chilevisión e Interferencia situaron el origen en la quebrada que baja al abanico de Álvaro del Portillo, La Plaza y General Blanche, donde la quebrada entra a un ducto bajo la urbanización. El escenario supone que el ducto no recibe flujo y asigna todo el aporte a la superficie desde ese punto. \emph{Cómo se representó.} Un hidrograma triangular de mezcla inyectado en la salida del cauce al abanico. \emph{Qué resultó.} En el modelo con solo el cauce principal (v01), el barro recorre General Blanche y Camino El Alba y llega cerca de Padre Hurtado, pero no reproduce la afectación informada en el sector sur: la Av.\ La Plaza Sur, la Clínica Universidad de los Andes y el campus quedan secos. En la prueba con observaciones reservadas alcanza los mismos 19 de 49 puntos que la hipótesis seleccionada. \emph{Estado.} La entrada por el cauce principal reproduce la rama norte en estos escenarios; por sí sola no reproduce la afectación informada en el sector sur.

\paragraph{H2: entrada superficial supuesta en el colector de Francisco Bulnes Correa.} \emph{Por qué se planteó.} Interferencia y BioBioChile recogieron versiones de un desborde donde el cauce entubado cruza Francisco Bulnes Correa, al inicio de Quebrada Honda, que figura entre las calles afectadas. \emph{Cómo se representó.} Una fuente en la embocadura del colector (e03), y escenarios hipotéticos de capacidad del colector como sumidero (v10, v11, v17, v18). \emph{Qué resultó.} Una salida solo en la embocadura no alcanza ninguno de los 49 puntos reservados: por sí sola no reproduce el conjunto de observaciones. Con el aporte supuesto desde el abanico, el modelo de ciudad encauza unos \SI{93000}{m^3} hacia la embocadura con picos cercanos a \SI{65}{m^3/s}. En esos escenarios, que no evalúan la obra real, el sumidero con capacidad supuesta de \SI{30}{m^3/s} alcanza ese límite y queda flujo simulado en Quebrada Honda. \emph{Estado.} Como entrada única, no reproduce las observaciones bajo los supuestos ensayados. Como punto de paso del flujo supuesto desde el abanico, reproduce flujo en Quebrada Honda; no establece por qué se afectó durante el evento.

\paragraph{H3: aporte repartido entre varios cauces.} \emph{Por qué se planteó.} Los sitios con afectación informada al sur (La Plaza Sur, el campus, la Clínica Universidad de los Andes) no son alcanzados en el modelo por el aporte del cauce principal. La cartografía empleada representa además un brazo sur que llega al campus y quebradas cortas que terminan en La Plaza Sur. \emph{Cómo se representó.} El mismo hidrograma, repartido en proporción al área de cada cuenca: \SI{79,5}{\percent} al cauce principal, \SI{18,8}{\percent} al brazo sur y \SI{1,7}{\percent} a las quebradas cortas. \emph{Qué resultó.} Sin lluvia sobre la ladera, la configuración con el brazo sur oficial (v21) obtiene un IEC de 68/100 y no alcanza la Clínica Universidad de los Andes. Con lluvia (H5, descrita a continuación), z02 obtiene 72/100 y fue seleccionada para la prueba reservada, desfavorable (19 de 49 puntos). \emph{Estado.} Es la base de todas las hipótesis posteriores, pero por sí sola deja sin barro gran parte de la zona afectada oficial.

\paragraph{H4: salida combinada.} \emph{Por qué se planteó.} Los primeros resultados mostraron que Quebrada Honda solo se mojaba si parte del flujo simulado salía en la embocadura. \emph{Cómo se representó.} Aportes en el abanico y en la embocadura, más los brazos sur (e01 y g03). \emph{Qué resultó.} En el cribado con \SI{150000}{m^3}, Quebrada Honda se moja en el \SIrange{54}{66}{\percent} de sus puntos, pero el frente queda a \SIrange{816}{1007}{m} de Padre Hurtado (sección~\ref{sec:res-calib}). \emph{Estado.} No cumple el criterio de alcance en esos ensayos.

\paragraph{H5: lluvia sobre la ladera.} \emph{Por qué se planteó.} Los primeros resultados mostraron que la Clínica Universidad de los Andes quedaba seca si solo llegaba barro por los cauces. \emph{Cómo se representó.} H4 más lluvia en el faldeo en el cribado (f02 y f03); después, cauces repartidos más lluvia con la ciudad completa (v02 y z02). \emph{Qué resultó.} z02 obtiene un IEC de 72 y alcanza 19 de 49 equipamientos en la prueba reservada, desfavorable. Sin la lluvia supuesta sobre la ladera (v21), el modelo no alcanza la Clínica Universidad de los Andes y el IEC baja a 68. \emph{Estado.} La lluvia supuesta sobre el faldeo contribuye a reproducir la afectación informada en el sector sur; la prueba reservada sigue siendo desfavorable. Su intensidad se revisa más abajo.

\subsection{Cuánto barro y de qué tipo}
\paragraph{Volumen.} \emph{Por qué importa.} En el modelo, el alcance depende del volumen supuesto: con \SI{150000}{m^3}, el frente se detiene a \SI{348}{m} de Padre Hurtado; con \SI{300000}{m^3} llega a menos de \SI{100}{m}. \emph{Cómo se acotó.} Con estimaciones independientes del modelo urbano (sección~\ref{sec:res-volumen}): un balance lluvia--escorrentía con saturación del suelo supuesta a partir del antecedente de que entre el 5 y el 7 de octubre habían caído unos \SI{95}{mm} en Río Molina \citep{cr2_2026} da una mediana de \SI{295000}{m^3}. \emph{Estado.} El volumen de entrada compatible con el alcance simulado coincide en orden de magnitud con el balance de lluvia bajo sus supuestos; no se ajustó para que coincidiera.

\paragraph{Concentración del barro.} \emph{Por qué importa.} Un barro más espeso se detiene antes; uno más líquido llega más lejos y se reparte. \emph{Qué resultó.} En el modelo, el barro espeso (Glenwood~1) se detiene a \SI{396}{m} de Padre Hurtado; el muy fluido (Aspen~2) obtiene 70/100. La cuadratura sobre la concentración (sección~\ref{sec:res-posterior}) da un área esperada de \SI{169}{ha} (\SIrange{144,5}{192,3}{ha}) y un IEC entre 71 y 73. \emph{Estado.} Dentro del rango plausible, la concentración cambia poco la comparación entre hipótesis.

\subsection{Cómo se representa la ciudad en el modelo}
\paragraph{Encauzamiento por las calles representadas.} \emph{Por qué se planteó.} El primer terreno tenía 1.103 edificios; el barro se repartía en superficie y no reproducía el metro informado en la Clínica UC Christus. \emph{Qué resultó.} Con las 42.018 huellas de edificios, el barro alcanza mayor profundidad en las calles: junto a la Clínica UC Christus sube de \SI{0,3}{m} a \SIrange{1,1}{1,2}{m} y la embocadura del colector recibe siete veces más barro. \emph{Estado.} Confirmado como efecto del modelo: sin los edificios, la reconstrucción no es comparable con lo observado.

\paragraph{Entrada de barro en las casas representadas.} \emph{Por qué se planteó.} Tras la prueba reservada, el diagnóstico mostró que el modelo dejaba seco el interior de la franja afectada: el barro corría por las calles pero no entraba a los lotes, mientras que en las imágenes consultadas se interpreta presencia de barro dentro de casas y patios. \emph{Cómo se representó.} Edificios de la zona afectada oficial con piso a \SI{0,3}{m} sobre el terreno y alta resistencia, que dejan entrar y retienen el barro (v26). \emph{Qué resultó.} La fracción del barro urbano dentro de la zona afectada oficial sube de 0,34 a 0,54 y el IEC verificado llega a 76. Con la mitad de la resistencia el resultado casi no cambia: lo que importa es que el barro entre a los lotes. La contrapartida es el espesor junto a la Clínica UC Christus, que baja a \SI{0,30}{m} porque las casas representadas aguas arriba almacenan barro en el modelo. Dejar sólidos los bloques grandes y porosas solo las casas de \SI{250}{m^2} o menos (d39) no recupera ese espesor (\SI{0,31}{m}): el efecto viene del almacenamiento aguas arriba, no del tamaño de los edificios.
\emph{Estado.} Mejora la coincidencia con la zona afectada oficial, pero es un ajuste posterior a la prueba reservada y no se ha validado con observaciones nuevas.

\subsection{El campus de la Universidad de los Andes y La Plaza Sur}
\paragraph{Aportes de la avenida y de la ladera ensayados.} \emph{Por qué se planteó.} En el video de Meganoticias se interpreta agua turbia que baja del bosque y llena un patio hundido del campus (espesor estimado visualmente: \SIrange{15}{35}{cm}), y El Dínamo recoge que Senapred nombró La Plaza Sur entre las calles más afectadas. El modelo de ciudad dejaba ese sector casi seco. \emph{Qué se ensayó.} Primero, un aporte superficial supuesto en Álvaro del Portillo hacia el campus: con ninguno de los fluidos ensayados, ni siquiera agua turbia, el flujo simulado llega al patio, porque el punto de aporte queda fuera de la cuenca local delimitada para el patio. Después, la escorrentía de la propia ladera y del cerro, con los rasgos que el relieve de \SI{30}{m} no resuelve: la calzada baja de La Plaza Sur, el paso del claro al patio y el camino de tierra del cerro. \emph{Qué resultó.} Con esos rasgos, el agua sigue La Plaza Sur hacia el norte, entra por el claro y llena el patio hasta \SIrange{0,12}{0,30}{m}, en el rango del video, con lluvias de \SI{15}{mm/h} a \SI{35}{mm/h}. \emph{Estado.} El recorrido simulado La Plaza Sur $\rightarrow$ claro $\rightarrow$ patio es compatible con lo observado y persiste con las intensidades de lluvia ensayadas. El aporte supuesto en Álvaro del Portillo no alcanza el patio bajo los supuestos ensayados. Con malla fina en el campus dentro de la simulación de ciudad y la lluvia medida (m36), el patio llega a \SI{0,17}{m} y La Plaza Sur se moja en el \SI{58}{\percent}: la diferencia entre escalas era de resolución.

\paragraph{Faldeo más líquido.} \emph{Por qué se planteó.} En el video citado de Meganoticias se interpreta agua turbia y no se distinguen escombros urbanos; esto motiva ensayar un flujo más diluido en el faldeo. \emph{Qué resultó.} Con una concentración menor en el faldeo, el campus se moja también con paso fino y el IEC verificado llega a 78. \emph{Estado.} Compatible: el escenario con un fluido más líquido en el sector sur reproduce mejor las observaciones de ajuste; no determina la composición real del flujo.

\paragraph{El \emph{bike park}.} \emph{Por qué se planteó.} En las imágenes consultadas, el autor identificó sobre el camino del cerro un \emph{bike park} con pistas de tierra desnuda; se ensayó su representación para examinar el recorrido calculado hacia el camino. \emph{Qué resultó.} Con y sin la pista principal tallada, los resultados en el camino, La Plaza Sur, el claro y el patio son idénticos (mrbc y mrbt): con este relieve, el agua ya escurría hacia el camino. Incluir toda la superficie del \emph{bike park} y de la ladera alta en la lluvia sí aumenta el agua (La Plaza Sur pasa de \SI{69}{\percent} a \SI{78}{\percent} mojada). \emph{Estado.} En los ensayos, tallar la pista no cambia el recorrido calculado; incluir esa superficie en la lluvia aplicada aumenta el aporte de agua del modelo. Su posible aporte de sedimento no puede evaluarse con este modelo, que no simula erosión.

\subsection{La hipótesis integrada y la lluvia}
\paragraph{Hipótesis integrada.} Reúne los cauces repartidos, la lluvia sobre el faldeo, las casas porosas, el faldeo líquido y los rasgos del campus. Con la lluvia supuesta de \SI{35}{mm/h} durante una hora obtiene 82/100, verificado con paso fino, y cumple los cuatro criterios originales. Con la lluvia medida en Río Molina y su forma real, un pulso con el máximo de \SI{10,9}{mm} en \SI{30}{min} (p34b, corregida), obtiene 81/100; el campus llega a \SI{0,49}{m}, La Plaza Sur se moja en el \SI{33}{\percent} de sus puntos y la superficie con más de \SI{5}{cm} es de \SI{222}{ha}. Con la misma lluvia repartida de forma uniforme en \SI{4}{h} (v33) baja a 79/100.
\emph{Estado.} Es la hipótesis más compatible con las observaciones de ajuste, pero se construyó después de la prueba reservada: validarla requiere observaciones independientes nuevas.

\paragraph{El obstáculo hipotético de árboles.} \emph{Por qué se planteó.} El autor marcó una línea de árboles que baja por la quebrada principal y una acumulación al final, poco antes de la salida al abanico. Se ensayó un obstáculo supuesto para examinar si el flujo calculado se desvía hacia el sur. \emph{Cómo se representó.} Un cauce tallado por la línea de árboles, con el flujo inyectado aguas arriba, y el mismo cauce con la acumulación representada como un obstáculo de \SI{3}{m} de altura (arc y art). \emph{Qué resultó.} En el modelo, el obstáculo supuesto no desvía el flujo hacia el sur: el sector de Álvaro del Portillo y el campus tiene la misma superficie mojada con y sin el obstáculo supuesto (\SI{20,2}{ha}). El obstáculo del modelo retiene y reparte barro junto a la salida al abanico, y el barro llega menos lejos hacia el poniente. El IEC casi no cambia (78 frente a 79). \emph{Estado.} Con este relieve, el obstáculo supuesto al final de la línea de árboles no reproduce la afectación informada en el sector sur.

\paragraph{El incendio.} \emph{Por qué se planteó.} Según el CR2, el incendio de diciembre de 2025 afectó la mayor parte de la cuenca. La literatura citada en el análisis de severidad describe posibles aumentos de escorrentía (suelo hidrofóbico, sin cobertura) y de material suelto disponible tras incendios. \emph{Cómo se cuantificó.} Con la severidad estimada por satélite (dNBR, sección~\ref{sec:res-incendio}): el \SI{82}{\percent} de las cuencas aportantes muestra señal de quema, pero solo un \SI{10}{\percent} con severidad moderada-alta o alta. Con esa severidad y los coeficientes supuestos, el balance del escenario sin incendio da un pulso de barro de alrededor del \SI{62}{\percent} del volumen con incendio (entre \SI{51}{\percent} y \SI{76}{\percent}), con el \SI{43}{\percent} de los sólidos. El modelo empírico Staley M1 da una probabilidad mediana de 0,28 con incendio frente a 0,07 sin él, como uso ilustrativo. \emph{Qué resultó en la ciudad.} Con el volumen sin incendio (f35b y f38b, corregidas), el barro moja \SIrange{179}{193}{ha}, frente a \SI{222}{ha} en p34b (\SIrange{13}{19}{\percent} menos), alcanza menos calles (0,70--0,73 frente a 0,77) y deja menos espesor en la Clínica UC Christus (\SIrange{0,19}{0,23}{m} frente a \SI{0,30}{m}). Sin embargo, el IEC no baja (80--82), porque premia la fracción del barro dentro de la zona afectada oficial, que sube cuando hay menos barro. \emph{Estado.} Con los coeficientes supuestos, el escenario con incendio calcula mayor volumen y, sobre todo, más sólidos y una mayor probabilidad empírica de flujo de detritos. El modelo urbano, con estas métricas, no permite afirmar que ese volumen adicional fuera necesario para reproducir lo observado.

\subsection{Aspectos no probados al cierre}
La lista completa figura en «Lo que permanece sin probar» (sección~\ref{sec:conclusiones}); aquí se destacan tres límites.
\begin{itemize}
\item \textbf{Validación independiente.} Todo lo posterior a la prueba reservada es ajuste. Su validación independiente permanece no probada: no se contrastó la hipótesis integrada con observaciones nuevas, como imágenes de radar posteriores al evento o marcas de nivel oficiales.
\item \textbf{La Clínica UC Christus.} Ninguna configuración reproduce a la vez el metro informado allí y la extensión dentro de la zona afectada oficial.
\item \textbf{Lo que el modelo no simula:} erosión y depósito durante el recorrido, troncos y vehículos arrastrados, el drenaje urbano, la infiltración dentro de la ciudad y los cambios de concentración del flujo.
\end{itemize}

\begin{table}[htbp]\centering\small
\caption{Resumen de las hipótesis de modelación: qué las motivó, qué se calculó y en qué estado quedan dentro de los ensayos. IEC: índice exploratorio de compatibilidad; QH: Quebrada Honda. Las compatibilidades se refieren al modelo; no identifican causas del evento.}\label{tab:recorrido}
\begin{adjustbox}{max width=\textwidth}\begin{tabular}{@{}>{\raggedright\arraybackslash}p{.17\textwidth}>{\raggedright\arraybackslash}p{.27\textwidth}>{\raggedright\arraybackslash}p{.30\textwidth}>{\raggedright\arraybackslash}p{.22\textwidth}@{}}\toprule
Hipótesis & Antecedente que la motivó & Resultado del modelo & Estado en el modelo\\\midrule
H1 cauce principal & Origen informado en el abanico (prensa citada en H1) & Reproduce la rama norte; no el sector sur & Insuficiente por sí sola\\
H2 colector Francisco Bulnes Correa & Desborde informado al inicio de Quebrada Honda (prensa citada en H2) & Como entrada única, 0 de 49 puntos; como paso del flujo, reproduce Quebrada Honda & Insuficiente como entrada única\\
H3 aporte repartido & Sector sur fuera del cauce principal & IEC 68 sin lluvia (v21); 72 con lluvia (z02, H5) & Base de las siguientes\\
H4 salida combinada & Quebrada Honda seca sin aporte en la embocadura & Cribado: QH \SIrange{54}{66}{\percent}; frente a \SIrange{816}{1007}{m} de Padre Hurtado & No cumple el alcance\\
H5 lluvia en la ladera & Clínica Universidad de los Andes seca en el modelo sin lluvia & IEC 72 (z02); sin lluvia, 68; prueba reservada 19/49 & Ajuste compatible; prueba desfavorable\\
Volumen & Alcance hasta Padre Hurtado & \SI{300000}{m^3} supuestos, del orden del balance con suelo saturado & Coherente\\
Edificios & Metro informado por Cadena 3 en la Clínica UC Christus & Las calles encauzan el barro & Confirmado en el modelo\\
Casas porosas & Interior de la zona afectada oficial seco en el modelo & IEC 76; Clínica UC Christus baja a \SI{0,30}{m} & Ajuste no validado\\
Campus por La Plaza Sur & Video de Meganoticias; Senapred según El Dínamo & Patio \SIrange{0,12}{0,30}{m} por el claro; m36: \SI{0,17}{m} y La Plaza Sur \SI{58}{\percent} & Compatible\\
Desborde de Álvaro del Portillo al campus & Ubicación supuesta del desborde & No llega al patio con los fluidos ensayados & No reproducido en los ensayos\\
\emph{Bike park} & Pistas desnudas sobre el camino & No cambia el recorrido; su superficie aporta agua & Recorrido sin efecto\\
Integrada & Todo lo anterior & IEC 79--82 según la lluvia & Mejor ajuste, sin validar\\
Obstáculo de árboles & Línea y acumulación identificadas en imágenes & No desvía el flujo al sur; retiene barro junto al abanico & No reproduce el sector sur\\
Incendio & Incendio informado por el CR2 nueve meses antes & Escenario sin incendio: aproximadamente \SI{62}{\percent} del volumen, \SI{43}{\percent} de los sólidos y aproximadamente 4 veces menos probabilidad; en la ciudad, \SIrange{13}{19}{\percent} menos superficie, IEC 80--82 frente a 81 & Sensibilidad a supuestos\\\bottomrule
\end{tabular}\end{adjustbox}
\end{table}
